\section*{Abstract}

This study was conducted with three main objectives: (1) the development of the EE-index, (2) the clarification of the occurrence characteristics of peculiar equatorial electrojet (EEJ) events, and (3) the development of a space weather database system.

First, a methodology for calculating the ``EE-index'', which quantitatively evaluates geomagnetic variations in the magnetic equatorial region, was established using data from the MAGDAS observation network operated by Kyushu University. In particular, by focusing on variations in the EUEL index, which represents localized disturbance components, it became possible to characterize the state of ionospheric currents in the equatorial region in detail.

Second, using the calculated EUEL index, an automatic detection method for ``peculiar EEJ'' events, which previously relied on visual inspection, was developed, and their occurrence characteristics were elucidated. Peculiar EEJ is a phenomenon in which the intensity difference of EUEL between the magnetic equator and low-latitude regions disappears. In this method, the difference in daily maximum EUEL values between the magnetic equator and low-latitude regions ($\Delta_{peak} EUEL$) was utilized, and peculiar EEJ events were quantitatively extracted by defining the lower 10\% of its distribution as a threshold. Furthermore, based on the correlation coefficient of the EUEL waveforms between the two sites, events were classified into an ``undeveloped type'' (correlation coefficient $\geq$ 0.6), in which the waveforms are similar, and a ``sudden type'' (correlation coefficient < 0.6), in which the waveforms are divergent.
Analysis of data from 2009 to 2020 resulted in the detection of 40 new events and the correction of misclassifications in previous visual inspections.
In this study, while an investigation covering the South American (Peru and Brazil) region was conducted following previous research, a more detailed analysis limited to the Brazilian region was also performed using data from the INTERMAGNET TTB station.
The results revealed that the undeveloped type tends to occur around the Southern Hemisphere winter solstice and at specific lunar ages. This suggests that seasonal variations in the Sq current system and lunar tidal effects are involved in the generation of peculiar EEJ events.

Third, a ``space weather database system'' was developed to widely provide the calculated EE-index and analysis results to the research community. This system features real-time calculation of indices from MAGDAS observation data and their visualization and publication on the web. This is expected to contribute to improving the accuracy of space weather monitoring and serve as a research infrastructure for related scientific fields.